\documentclass[aspectratio=169,spanish]{beamer}
%--------------------------------------------------
%  PAQUETES BÁSICOS
%--------------------------------------------------
\usepackage[utf8]{inputenc}
\usepackage[T1]{fontenc}
\usepackage{babel}
\usepackage{graphicx}
\usepackage{amsmath,amssymb}
\usepackage{tikz}
\usepackage{booktabs}

%--------------------------------------------------
%  TEMA Y COLORES
%--------------------------------------------------
\usetheme{Madrid}
\usecolortheme{beaver}

%--------------------------------------------------
%  DATOS DE LA PRESENTACIÓN
%--------------------------------------------------
\title[Transporte ZMG y Teoría de Grafos]{El transporte urbano de la Zona Metropolitana de Guadalajara\\desde la perspectiva de la teoría de grafos}
\author[F. Betancourt]{Fernando Betancourt Moreno\\\small Director: Dr. Abel Palafox González}
\institute[CUCEI -- UdeG]{Universidad de Guadalajara\\Centro Universitario de Ciencias Exactas e Ingenierías}
\date[Feb.~2025]{Febrero de 2025}

%--------------------------------------------------
%  DIAPOSITIVA DE SECCIÓN
%--------------------------------------------------
\AtBeginSection[]{%
  \begin{frame}[plain]
    \centering\Huge{\insertsection}
  \end{frame}}

%--------------------------------------------------
%  INICIO DEL DOCUMENTO
%--------------------------------------------------
\begin{document}

%====================
%  PORTADA & ÍNDICE
%====================
\begin{frame}[plain]
  \titlepage
\end{frame}

\begin{frame}{Contenido}
  \tableofcontents
\end{frame}

%====================
%  INTRODUCCIÓN
%====================
\section{Introducción}

\begin{frame}{Contexto y problemática}
  \begin{itemize}
    \item Expansión urbana acelerada $\rightarrow$ crecimiento desordenado en la ZMG.
    \item Red de transporte actual: saturación pico, cobertura desigual, baja intermodalidad.
    \item Necesidad de diagnósticos cuantitativos para guiar políticas de movilidad sostenible (IMEPLAN, 2016).
  \end{itemize}
\end{frame}

\begin{frame}{Relevancia y alcance}
  \begin{itemize}
    \item Transporte eficiente = calidad de vida + competitividad regional.
    \item La \textbf{teoría de grafos} permite:
      \begin{itemize}
        \item Detectar cuellos de botella y hubs críticos.
        \item Medir accesibilidad y rodeo.
        \item Simular fallas y proponer reforzamientos.
      \end{itemize}
    \item Resultados extrapolables a otras metrópolis latinoamericanas.
  \end{itemize}
\end{frame}

%====================
%  OBJETIVOS E HIPÓTESIS
%====================
\section{Objetivos e hipótesis}

\begin{frame}{Pregunta de investigación}
  \begin{block}{Pregunta central}
    \centering
    \emph{¿Cómo puede la teoría de grafos proporcionar un marco metodológico para analizar y optimizar la conectividad del sistema de transporte público en la ZMG?}
  \end{block}
\end{frame}

\begin{frame}{Objetivos}
  \textbf{General}
  \begin{itemize}
    \item Evaluar la conectividad, accesibilidad y vulnerabilidad de la red de transporte público.
  \end{itemize}
  \vspace{0.4cm}
  \textbf{Específicos}
  \begin{enumerate}
    \item Identificar nodos y rutas con mayor centralidad.
    \item Analizar eficiencia temporal y cobertura espacial.
    \item Evaluar resiliencia ante fallas en hubs críticos.
    \item Contrastar con casos de referencia (Sydney, Resistencia).
  \end{enumerate}
\end{frame}

\begin{frame}{Hipótesis}
  Modelar la red con métricas gráficas (betweenness, spectral gap, índice $\alpha$) permitirá sugerir intervenciones que reduzcan tiempos de viaje y aumenten robustez frente al diseño actual.
\end{frame}

%====================
%  MARCO TEÓRICO
%====================
\section{Marco teórico}

\begin{frame}{Teoría de grafos \& transporte}
  \begin{itemize}
    \item Nodo $\leftrightarrow$ parada o estación.
    \item Arista $\leftrightarrow$ tramo (ponderado por tiempo, distancia o tarifa).
    \item Indicadores globales: $\beta$, $\gamma$, diámetro, brecha espectral.
    \item Indicadores nodales: grado, cercanía, betweenness.
  \end{itemize}
\end{frame}

\begin{frame}{Modelos y datos de soporte}
  \begin{columns}
    \column{0.48\textwidth}
      \textbf{GTFS}
      \begin{itemize}
        \item Estándar abierto: rutas, horarios, tarifas.
        \item Base para simulaciones en \texttt{OpenTripPlanner}.
      \end{itemize}
    \column{0.48\textwidth}
      \textbf{Algoritmos de ruta mínima}
      \begin{itemize}
        \item Dijkstra (fuente única).
        \item Floyd–Warshall (todo par).
      \end{itemize}
  \end{columns}
\end{frame}

%====================
%  METODOLOGÍA
%====================
\section{Metodología}

\begin{frame}{Pasos del estudio}
  \begin{enumerate}
    \item Levantamiento GTFS y EOD~2023.
    \item Construcción del grafo dirigido ponderado.
    \item Cálculo de métricas (\texttt{networkx}/\texttt{igraph}).
    \item Simulación de fallas y escenarios de expansión.
    \item Validación con tiempos de viaje de \texttt{OTP}.
  \end{enumerate}
\end{frame}

\begin{frame}{Métricas clave}
  \begin{block}{Nivel de red}
    $\beta$ (conectividad), $\alpha$ (redundancia), $DI$ (rodeos), $\Delta$ (brecha espectral).
  \end{block}
  \vspace{0.4cm}
  \begin{block}{Nivel de nodo}
    Grado, centralidad de cercanía, betweenness, eccentricidad.
  \end{block}
\end{frame}

%====================
%  CASOS COMPARATIVOS
%====================
\section{Casos comparativos}

\begin{frame}{Sydney (Zhao et al., 2024)}
  \begin{itemize}
    \item 3,799 paradas; 76 rutas de autobús + metro.
    \item Modelo gravitacional OD con costos topológicos.
    \item Ponderación entrópica redujo RMSE 4\,\%.
  \end{itemize}
\end{frame}

\begin{frame}{Resistencia (Cardozo et al., 2009)}
  \begin{itemize}
    \item 48 nodos – 61 arcos.
    \item $\beta = 1.27$, $\alpha = 0.154$.
    \item Gradiente centro–periferia pronunciado.
  \end{itemize}
\end{frame}

\begin{frame}{Guadalajara (Barraza \& Estrada, 2021)}
  \begin{columns}
    \column{0.52\textwidth}
      \begin{itemize}
        \item 2,033 nodos, 4,704 aristas.
        \item Tiempo medio de viaje: 70 min.
        \item $DI = 0.63$ (trayectos relativamente rectos).
      \end{itemize}
    \column{0.48\textwidth}
      \begin{itemize}
        \item Tres hubs concentran $>40\%$ del flujo.
        \item Vulnerabilidad alta a fallas localizadas.
      \end{itemize}
  \end{columns}
\end{frame}

%====================
%  HALLAZGOS DE LA LITERATURA
%====================
\section{Hallazgos de la literatura}

\begin{frame}{Puntos clave – Barraza \& Estrada (2021)}
  \begin{itemize}
    \item Índice $\beta = 1.05$ \textrightarrow{} topología casi arborescente.
    \item Super‑hubs: Centro Histórico, Calzada Lázaro Cárdenas, Periférico Sur.
    \item Remoción del top 10\% de nodos incrementa la longitud de camino promedio 91\% (APL = 69 paradas).
  \end{itemize}
  \medskip
  {\small\itshape Fuente: Barraza \& Estrada, 2021.}
\end{frame}

%====================
%  LIMITACIONES Y FUTURO TRABAJO
%====================
\section{Limitaciones y futuro}

\begin{frame}{Limitaciones actuales}
  \begin{itemize}
    \item Datos OD agregados limitan la resolución espacial de los flujos.
    \item Supuestos de homogeneidad vehicular y puntualidad pueden subestimar la variabilidad real.
    \item No se incluyen factores socio-económicos ni topográficos que condicionan la demanda.
  \end{itemize}
\end{frame}

%====================
%  CONCLUSIONES
%====================
\section{Conclusiones}

\begin{frame}{Conclusiones}
  \begin{itemize}
    \item La teoría de grafos ofrece un diagnóstico robusto y replicable de la red ZMG.
    \item Tres super-hubs concentran el flujo, lo que implica alta vulnerabilidad estructural.
    \item Intervenciones dirigidas (más conexiones periféricas y diagonales) podrían reducir tiempos de viaje en >15 \% y aumentar resiliencia en >20 \%.
  \end{itemize}
\end{frame}

%====================
%  CIERRE
%====================
\section*{Preguntas}

\begin{frame}[plain]
  \centering
  \Huge\bfseries ¿Preguntas?
\end{frame}

\end{document}

